% ============================================================================
% Fichier  : front/01_avant_propos.tex
% Titre    : Avant-propos
% Auteur   : MINKA MI NGUIDJOÏ Thierry Emmanuel
% Version  : 1.0 -- novembre 2026
% ============================================================================

\chapter*{Avant-propos}
\addcontentsline{toc}{chapter}{Avant-propos}
\markboth{Avant-propos}{Avant-propos}

\epigraph{« Le risque se gère avant, la crise pendant,\\
et l'apprentissage après. »}{Principe du cours}

Ce cours s'adresse aux élèves ingénieurs de quatrième année du Département
de Génie Informatique de l'École Nationale Supérieure Polytechnique de
Yaoundé. Il suppose acquis le cours de sécurité informatique (SEC3266) et
prolonge l'audit des systèmes d'information (INF3226)~: il ne réexplique pas
les mesures techniques, il apprend à décider à leur sujet, à vérifier
qu'elles tiennent, et à diriger l'organisation quand elles ont cédé.

Il repose sur deux convictions. La première~: un risque n'est pas une
couleur dans une matrice, mais une incertitude sur des pertes, que l'on
peut estimer, discuter et réduire. La seconde~: une crise ne se gère pas
avec un plan, mais avec une organisation entraînée, qui sait qui décide,
qui informe et qui consigne. Entre les deux, l'audit vérifie que ce que
l'on croit en place l'est vraiment, et la continuité garantit que
l'essentiel reprend dans les délais convenus.

\bigskip

\textbf{Ce que ce cours n'est pas.} Il n'est pas une préparation à une
certification, ni un recueil de modèles à remplir, ni un cours de
communication institutionnelle. Les méthodes y sont présentées avec leurs
limites, et les chiffres avec leur marge.

\bigskip

\textbf{Ce que ce cours est.} Un traité de vingt-six chapitres, en six
parties précédées d'un prologue. La Partie~I pose les fondements~: risque
et incertitude, perception et biais, gouvernance, cadres de référence,
obligations juridiques. La Partie~II apprend à analyser~: identifier,
coter, conduire EBIOS Risk Manager, quantifier avec FAIR, simuler par
Monte-Carlo, traiter et rendre compte. La Partie~III apprend à auditer le
dispositif selon l'ISO~19011. La Partie~IV traite des incidents et de la
continuité. La Partie~V est consacrée à la gestion de crise~: cellule,
décision, communication, exercices. La Partie~VI apprend à tirer les leçons
et à mesurer la résilience, y compris ce que l'on en ignore.

\bigskip

\textbf{Un fil rouge.} Le cours suit l'entreprise fictive LiZbiZ, société
de logistique aux quatre centres de données, frappée par un rançongiciel
une nuit de mars (Opération GAROUA)~: le cours de sécurité en a étudié les
causes techniques~; celui-ci étudie la décision, avant, pendant et après.
L'Opération LIZBIZ suit le programme de gestion des risques et l'audit qui
ont suivi~; l'Opération WOURI, la continuité de la BANQUE SANAGA~;
l'Opération MOUNGO, l'inondation d'un entrepôt de Douala. Les noms, les
dates et les montants sont fictifs~; les méthodes sont exactes, et les
textes de loi sont cités d'après leur version officielle.

\bigskip

\textbf{Comment travailler.} Lisez chaque chapitre avec sa fiche
d'auto-évaluation~; faites les exercices avant de lire leurs corrigés~;
pratiquez dans les labs de NDAP SUKULU. Prenez l'habitude que le cours
cherche à transmettre~: séparer ce que l'on a constaté de ce que l'on en
pense, donner une marge avec chaque chiffre, et dire ce que l'on ne sait
pas.

\bigskip

\begin{flushright}
    \IfFileExists{../../Logos/signature_avant_propos.png}{\includegraphics[height=1.6cm]{../../Logos/signature_avant_propos.png}}{\textit{MINKA MI NGUIDJOÏ Thierry Emmanuel}}\\
    \textit{Yaoundé, novembre 2026}
\end{flushright}
